\subsection{Materials: Electrical Behavior and Suitability for Phantom Construction}

The phantom head design relies on a bilayer architecture composed of an external conductive
shell and an internal conductive volume. These two layers are subject to fundamentally
different material constraints. The outer shell is fabricated using a commercially
available conductive polymer, selected for accessibility, cost, and mechanical
printability. While alternative formulations or custom composites could, in principle, be
developed, the electrical properties of the shell material are largely fixed by the
commercial product and the 3D printing process. As such, its characterization is primarily
based on manufacturer data and literature, rather than on experimental tunability.

In contrast, the inner conductive layer is based on a saline--gelatine medium, which is
readily available, inexpensive, and highly customizable through formulation parameters.
For this reason, the experimental characterization efforts focused on the inner layer, with
the goal of assessing its electrical behavior, tunability, and reproducibility in the
context of phantom construction. Rather than aiming at a precise metrological determination
of conductivity values, the experiments were designed to qualitatively evaluate how the
electrical response of the material evolves with ionic concentration, and whether these
properties can be reliably reproduced across fabrication batches.

The characterization was performed using three formulations with NaCl concentrations of
0\%, 2.5\%, and 5\%. These values were deliberately chosen to span a range of electrical
behaviors, from predominantly resistive to increasingly reactive, and to ensure that
capacitive effects would be observable within the resolution limits of the available
measurement equipment. As the salt concentration increases, the magnitude of the impedance
decreases across the explored frequency range, reflecting the expected increase in ionic
conductivity. In parallel, the phase response progressively departs from purely resistive
behavior, with higher concentrations exhibiting increasingly negative phase angles at low
frequencies.

The emergence of this reactive component is attributed to interfacial polarization effects
at the electrode--gelatine interface. At higher ionic strengths, the increased density of
mobile charge carriers promotes the formation of an electrical double layer, which enhances
the effective capacitance of the system, particularly in the low-frequency regime relevant
to EEG measurements. These observations are consistent with prior studies on gelatine-based
phantoms. In particular, Owda~\emph{et al.}~\cite{Owda2020} reported analogous trends in
impedance magnitude and phase response as a function of salt concentration, including the
appearance of capacitive effects at higher salinities. The agreement in qualitative behavior
supports the interpretation that the observed phase shifts arise from intrinsic
material--electrode interactions rather than artifacts of the specific experimental setup.

While the higher concentrations explored in this study are not representative of the
intended operating regime of the final phantom, they serve to clearly delineate the
transition from predominantly Ohmic to increasingly reactive behavior. From a design
perspective, this regime separation is informative, as it demonstrates that the electrical
properties of the gelatine medium are systematically adjustable through ionic concentration.
For the final phantom implementation, lower concentrations will be selected to ensure
operation within the linear and Ohmic regime required by the electroquasistatic modeling
framework.

Beyond tunability, a key outcome of the impedance analysis is the high degree of
repeatability observed between batches prepared with identical formulations. Both impedance
magnitude and phase responses exhibited negligible variation across repeated fabrications.
This batch-to-batch consistency is a critical requirement for phantom construction, as it
ensures that multiple physical realizations can be produced with stable and reproducible
electrical properties, enabling meaningful comparison across experiments and long-term use
of the phantom platform.

The qualitative signal localization experiments further support the suitability of the
gelatine medium as a volume conductor. The measured potentials varied coherently with
electrode position and relative distance to multiple internal sources, with measurement
locations closer to a given source exhibiting higher signal amplitudes at the corresponding
excitation frequency. Although no quantitative localization accuracy is claimed, these
results confirm that the medium supports spatially structured potential distributions and
allows the superposition and differentiation of simultaneous sources—an essential
requirement for realistic phantom operation.
